\uuid{th3l}
\titre{Variation de la constante/prolongement par continuité}
\auteur{hanani ; mzari}
\organisation{exo7}
\licence{CC BY-NC-SA 4.0}
\datecreate{2020-01-20}
\niveau{L1}
\module{Analyse}
\chapitre{Equation différentielle}
\sousChapitre{Résolution d'équation différentielle du premier ordre}
\difficulte{3}
\theme{Equations du premier ordre, variation de la constante, prolongement par continuité}
\exo7qcmid{lille-2.equations-differentielles.016}

\contenu{
\begin{question}
\qclassification{225.01}
\qtags{motcle=variation de la constante/prolongement par continuité}
On considère l'équation différentielle $(E)$ : $y'-y\tan x=1$ sur $]-\pi/2,\pi /2[$. Parmi les affirmations suivantes, cocher celles qui sont vraies :
\begin{answers}
\good{La solution générale de l'équation homogène est $\displaystyle y=\frac{k}{\cos (x)}$, $k\in \Rr$.}
\good{Si $\displaystyle y=\frac{k(x)}{\cos (x)}$ est une solution de $(E)$, alors $k'(x)=\cos (x)$.}
\bad{La solution générale de $(E)$ est $\displaystyle y=\frac{k}{\cos (x)}+\sin x$, $k\in \Rr$.}
\bad{$(E)$ possède une solution qui se prolonge par continuité en $-\pi/2$ et en $\pi /2$.}
\end{answers}
\begin{explanations}
Une primitive de $-\tan (x)$ est $\ln |\cos x|$, donc la solution générale de l'équation homogène est : $\displaystyle y(x)=\frac{k}{\cos (x)}$. La variation de la constante montre que $\displaystyle y_0=\frac{\sin (x)}{\cos (x)}$ est une solution particulière. Enfin, une solution $\displaystyle y=\frac{k+\sin (x)}{\cos (x)}$ de $(E)$ se prolonge par continuité en $-\pi/2$ si et seulement si $k=1$ et se prolonge par continuité en $\pi/2$ si et seulement si $k=-1$. Aucune solution de $(E)$ ne se prolonge par continuité en $-\pi/2$ et en $\pi /2$.
\end{explanations}
\end{question}
}
